%!TEX root = ./../../thesis-trier.tex

Here, in \cref{constrsync:prop2:SI-binary-sync-in-L,syncmax:SI-aut:prop:sync-L-complete}, 
we prove \LOGSPACE-completeness of deciding each of the properties 
of synchronizability, complete reachability, 
complete distinguishability, or sync-maximality
for binary SI-automata as input. \Cref{constrsync:prop2:SI-binary-sync-in-L}
is an easy consequence of the following lemma (we note that it was essentially
proved by Martyugin~\cite[Prop.~6.2]{Mar2009b}
that the problem from \cref{constrsync:prop:SI-binary-sync-in-L} 
can be solved in $O(n)$).

\begin{lemma}
	\label{constrsync:prop:SI-binary-sync-in-L}
	For $n \ge 4$, which case  from~\cref{prop:structure_simp_idem_binary} 
	applies and the computation of the respective parameters $s,t \in Q$, $a,b \in \Sigma$, 
	i.e., which letter is the simple idempotent $a$ and which cyclically permutes the states as~$b$,
	and $0 < d < n$ can be done in deterministic log-space.
\end{lemma}
\begin{proof}
	We denote by $\mathcal A_t$ the synchronizing 
	binary SI-automaton from~\cref{prelim:fig:binarySI-sync-trap} with a trap state
	and by $\mathcal A_c$ the other automaton from~\cref{prelim:fig:binarySI-sync-SCC}
	where $b$ permutes the states,~$a$ is a simple idempotent,
	and $s,t \in Q$ and $0 < d < n$
	are the parameters as in \cref{prop:structure_simp_idem_binary} w.r.t.~$\mathcal A_t$ and $\mathcal A_c$
	for $n \ge 4$. 
	We show a scheme to test whether $\mathcal A$
	equals~$\mathcal A_c$ or~$\mathcal A_t$, which gives the claim by \cref{prop:structure_simp_idem_binary}.
	By iterating through the states,
	we search for $t \in Q$
	with $\delta^{-1}(t, a) = \{s,t\}$
	for some $a \in \Sigma$ and $s \in Q$.
	%	with two incoming transitions
	%	and $\applytrans{\delta}{t}{a} = t$ for some $a \in \Sigma$.
	% In case we do not find such a state, then $\mathcal A$ is not synchronizing.
	% Otherwise, let $b \in \Sigma \setminus \{a\}$.
	Let $b \in \Sigma\setminus \{a\}$.
	Then, $\mathcal A$ equals~$\mathcal A_c$ or $\mathcal A_t$
	precisely in the following two cases.
	\[
	\mathcal A
	=
	\begin{cases}
		\mathcal A_c & \text{if } \delta(s, b^n) = s, \delta(s, b^d) = t,   \forall i \in \{1,\ldots,n-1\} \colon \delta(s, b^i) \ne s, \\
		\mathcal A_t & \text{if } \delta(s, b^{n-1}) = s, \delta(t, b) = t, \forall i \in \{1,\ldots,n-1\} \colon \delta(s, b^i) \ne s, 
	\end{cases}
	\]
	for some $d \in \{1,\ldots,n-1\}$ coprime to $n$.
	By iterating through the states $\delta(s, b^i)$
	and by \cref{prelim:prop:coprime_in_L} each case can be tested
	in deterministic log-space.
	%	
	%%	\begin{enumerate}
		%\ifthenelse{\equal{\thing}{CIAA}}
		%{
			%	\item For $\applytrans{\delta}{t}{b} \ne t$ 
			%	iterate through $\applytrans{\delta}{t}{b^i}$
			%	for $i \in \{1,\ldots,n\}$
			%	and test whether $\delta(t, b^n) = t$, $\delta(t, b^i) \ne t$ for $i \in \{2,\ldots,n-1\}$,
			%	$\delta(t, b^{n-d}a) = t$
			%	for some $d \in \{1,\ldots,n-1\}$
			%	coprime to $n$
			%	and $\delta(t, b^{i}a) = \delta(t, b^i)$
			%	for $i \in \{1,\ldots,n-1\} \setminus \{d\}$.
			%	This can be done in deterministic log-space and if all these tests are positive,
			%	then $\mathcal A = \mathcal A_c$.
			%}
		%{
			%\begin{comment} 
			%	\item For $\applytrans{\delta}{t}{b} \ne t$ (which excludes $\mathcal A_t$) iterate through $\applytrans{\delta}{t}{b^i}$
			%    for $i \in \{1,\ldots,n\}$
			%    and test whether $\delta(t, b^i) \ne t$ for $i \in \{2,\ldots,n-1\}$, $\delta(t, b^n) = t$, and
			%    $\delta(t, b^{n-d}a) = t$
			%    for some $d \in \{1,\ldots,n-1\}$
			%    coprime to $n$.
			%    %and $\delta(t, b^{i}a) = \delta(t, b^i)$
			%    %for $i \in \{1,\ldots,n-1\} \setminus \{d\}$.
			%    This can be done in deterministic log-space
			%    by \cref{prelim:prop:coprime_in_L}. In case all these tests are positive,
			%    we have $\mathcal A = \mathcal A_c$.
			%\end{comment} 
			%}    
		%\begin{comment} 
		%    \item For $\applytrans{\delta}{t}{b} = t$, test whether $\mathcal A = \mathcal A_t$ 
		%    by iterating through  $\delta(s, b^i)$, $i \in \{1,\ldots,n-1\}$,
		%    and checking $\delta(s, b^i) \ne s$
		%    for $i \ne n - 1$ and 
		%    $\delta(s, b^{n-1}) = s$.~\qedhere
		%%	\end{enumerate}	
	%\end{comment}        
\end{proof}

\begin{proposition} \label{constrsync:prop2:SI-binary-sync-in-L}
	Deciding whether a given binary $n$-state SI-automaton $\semiautomat{A}{Q}{\Sigma}{\delta}$ is synchronizing
	is in \LOGSPACE. 
\end{proposition}
\begin{proof}
	A binary SI-automaton is synchronizing
	if and only if it has one of the two forms
	from \cref{prop:structure_simp_idem_binary}, 
	and this can be tested in deterministic log-space
	by \cref{constrsync:prop:SI-binary-sync-in-L}.
\end{proof}

\begin{theorem} \label{syncmax:SI-aut:prop:sync-L-complete}
	For a binary SI-automaton it is \LOGSPACE-complete
	to decide each of the following properties individually: synchronizability,
	complete reachability, complete distinguishability, sync-maximality.	
\end{theorem}
\begin{proof}
	We can assume the input automaton has at least three states.
	Containment in \LOGSPACE
	follows from \cref{constrsync:prop:SI-binary-sync-in-L} 
	for synchronizability.
	By \cref{sync_max:prop:binary-SI}, 
	for a binary SI-automaton
	the properties of complete reachability, complete distinguishability, and sync-maximality
	are each equivalent to the fact that the automaton
	equals the one from \cref{prelim:fig:binarySI-sync-SCC},
	which can be decided in \LOGSPACE 
	by \cref{constrsync:prop:SI-binary-sync-in-L}.
	%	Next, note that for binary SI-automata with at least
	%	three states, the automaton in \cref{prelim:fig:binarySI-sync-SCC}
	%	is the only binary SI-automaton that is completely reachable,
	%	the only binary SI-automaton that is completely distinguishable,
	%	and the only binary SI-automaton that is sync-maximal,
	%	as the other synchronizing automaton is not strongly
	%	connected and two sets $A$, $A \cup \{t\}$
	%	with trap state $t$ are indistinguishable.
	%	So, we only have to test whether the input automaton
	%	equals this automaton, which can be done in \LOGSPACE
	%	by \cref{constrsync:prop:SI-binary-sync-in-L}.
	For \LOGSPACE-hardness, we give a reduction from \SCP. Let $f \colon [n] \to [n]$
	be a permutation with $n \ge 3$.
	We can assume $f \ne \operatorname{id}_{[n]}$
	with $f(p) \ne p$ for $p \in [n]$,
	which can be determined by an $\AC^0$ circuit.
	Set $\semiautomat{A}{[n]}{\{a,b\}}{\delta}$
	with $\delta(q, b) = f(q)$ and
	\[
	\delta(q, a) = \begin{cases}
		f(p) & \text{if } q = p, \\
		q    & \text{if } q \ne p.
	\end{cases}
	\]
	for $q \in [n]$. 
	Observe that $\mathcal A$ cannot equal the automaton from \cref{prelim:fig:binarySI-sync-trap}, as $a$ moves a point in a cycle of $f$
	to another point in the same cycle. Now, the following are equivalent:
	\begin{enumerate}
		\item $f$ cyclically permutes $[n]$,
		
		\item $\mathcal A$ is synchronizing,
		
		\item $\mathcal A$ has the form
		from \crefitem{sync_max:prop:binary-SI}{sync_max:cor:item:form},
		where $d = 1$ in this case.		
	\end{enumerate}
	So, by \cref{sync_max:prop:binary-SI},
	$f$ cyclically permuting $[n]$
	is also equivalent to 
	complete
	reachability and to sync-maximality.  
	Hence, the construction gives a valid reduction
	for all these properties.		
\end{proof}
